\subsubsection{CLAHE} \label{CLAHE}

Contrast Limited Adaptive Histogram Equalization (CLAHE) is an image enhancement technique that improves the contrast and brightness of an image. This technique is widely used in image processing and computer vision tasks to improve the visual quality of images and make it easier to detect objects and features, example of CLAHE is shown in Figure \ref{fig:CLAHE}.

\begin{figure}[!h]
	 \centerline{\includegraphics[width=0.8\textwidth]{src/augmentations/CLAEH.png}}
	 {\caption{Examples of CLAHE with contrast limit between 1 and 4 and title grid with size of 8} \label{fig:CLAHE}}
\end{figure}

CLAHE works by dividing the image into small rectangular regions called tiles and then equalizing the histogram of each tile separately. This approach ensures that the brightness and contrast of each tile are optimized independently, which can be beneficial in cases where the brightness or contrast varies significantly across the image.

One advantage of CLAHE is that it is an adaptive technique, meaning that it can adapt to the local characteristics of the image. This allows it to provide better results than traditional histogram equalization, which applies the same transformation to the entire image. CLAHE also allows for better visualization of fine details and edges in images, which can be especially useful in medical imaging applications or aerial images of marine wildlife where subtle changes in texture or contrast can be important for detecting features or anomalies.

\subsubsection{Sharpen} \label{Sharpen}

Sharpen augmentation is a technique used in image processing and computer vision that enhances the edges and fine details in an image. The technique involves applying a filter to the image that increases the contrast between adjacent pixels, making the edges more distinct and prominent.

Sharpen augmentation can be particularly advantageous in the context of aerial images for a number of reasons. Aerial images of the marine environment, for instance, often contain low-contrast regions such as large bodies of water or sky that can reduce the sharpness and clarity of the image. Sharpen augmentation can help to increase the contrast between adjacent pixels, making edges and fine details more visible and distinct. This can improve the visual quality of the image and make it easier to detect and recognize objects, including marine wildlife, in the image. 

The alpha parameter in sharpen augmentation is a factor that controls the strength of the filter. A higher value of alpha will result in a stronger sharpening effect, while a lower value will result in a weaker effect.  The lightness parameter in sharpen augmentation is a factor that controls the amount of lightness or brightness in the image. A higher value of lightness will increase the brightness of the image, while a lower value will decrease it. Examples of augmented images are represented in Figure \ref{fig:Sharpen}

\begin{figure}[!h]
	 \centerline{\includegraphics[width=0.8\textwidth]{src/augmentations/Sharpen.png}}
	 {\caption{Examples of Sharpen with alpha value between 0.2 and 0.5, lightness value between 0.5 and 1.0} \label{fig:Sharpen}}
\end{figure}

\subsubsection{Channel Shuffling}\label{ChannelShuffles}
Channel Shuffling augmentation is an image transformation technique that can be used to generate a diverse set of training images for machine learning models. In channel shuffles, the color channels of an image are randomly permuted, which can introduce a significant amount of variation to the image data.

The advantage of Channel shuffling is its simplicity in both understanding and implementation. However, the variation of the transformation is limited to the permutation of the channels, which is only six. This limitation restricts the ability to enhance the robustness of the data. Moreover, one color channel may be more important than others. For example, if blue is an important feature to recognize the ocean, shuffling the channels could pose an obstacle for the model. Figure \ref{fig:ChannelShuffles} show the effect of Channel Shuffling on the image.
\begin{figure}[!h]
	 \centerline{\includegraphics[width=0.8\textwidth]{src/augmentations/ChannelShuffle.png}}
	 {\caption{Examples of Channel Shuffling} \label{fig:ChannelShuffles}}
\end{figure}

\subsubsection{Color Jitter}\label{ColorJitter}
Color Jitter is an image augmentation technique that introduces random color variation to an image by perturbing the color channels of the image. The goal of this technique is to increase the diversity of the training data-set for machine learning models, which can help improve the model's ability to generalize to new, unseen data.

Color Jitter works by randomly perturbing the hue, saturation, brightness, and contrast of an image. These perturbations are generated by adding small amounts of noise or by scaling the color values in channels. By applying these perturbations, Color Jitter creates new images that are visually distinct from the original image, which can help improve the model's ability to detect and recognize objects with varying color properties.

\begin{figure}[!h]
	 \centerline{\includegraphics[width=0.8\textwidth]{src/augmentations/ColorJitter.png}}
	 {\caption{Examples of Color Jitter with the random changes in brightness, contrast, saturation, hue between 0 and 0.2} \label{fig:ColorJitter}}
\end{figure}

Similar but more sophisicated than \hyperlink[\ref{ChannelShuffles}]{Channel Shuffling}, Color Jitter can be advantageous for aerial images because it introduces random color variation to the images. For example, when detecting marine wildlife in aerial images, variations in the lighting and coloring of the water and surrounding environment can make it challenging to detect and differentiate the animals. By using Color Jitter to generate a diverse set of training images with random color perturbations, the machine learning model can learn to detect marine wildlife in a wider range of visual conditions. Additionally, standardizing the color distribution across images in the data-set can help to reduce the impact of color biases in the data-set, which can lead to improved accuracy and performance of the machine learning model, especially in classes with small samples. Figure \ref{fig:ColorJitter} illustrates the effect of the transformations.


\subsubsection{Solarize} \label{Solarize}
Solarize can be advantageous for aerial images, as it introduces a striking visual effect by inverting the colors of an image that exceed a certain threshold.

In the context of aerial images of marine wildlife, solarize can be useful to create a diverse set of training images for machine learning models. By inverting the colors of the images, solarize can create visually distinct images that can improve the model's ability to detect and recognize objects with varying color properties. Similar to Color Jitter \ref{ColorJitter}, Solarize helps to reduce the impact of color biases in the dataset. For example, if the majority of images of a particular marine species are predominantly of a particular color, the machine learning model may learn to associate that color with that class, leading to reduced accuracy and performance when detecting that species in images with different colors. By using solarize to invert the colors of these images, the model can be trained to recognize the species regardless of the color of the image.
\begin{figure}[!]
	 \centerline{\includegraphics[width=0.8\textwidth]{src/augmentations/Solarize.png}}
	 {\caption{Examples of Solarize with threshold of 128} \label{fig:cleah}}
\end{figure}

% \subsubsection{Normalize} \label{Normalize}

% Normalization is an image augmentation technique that involves scaling the pixel values of an image so that they fall within a specified range. This can be done by subtracting the mean pixel value of the image and dividing the result by the standard deviation of the pixel values. The concrete formula is:

% \begin{equation}
% img = \frac{(img - mean \cdot max\_pixel\_value)}{(std \cdot max\_pixel\_value)}
% \end{equation}

% As shown in in Figure \ref{fig:Normalize}, by scaling the pixel values to a standardized range, the machine learning model can better distinguish the features and patterns in the image that are relevant to the task, while reducing the impact of irrelevant variations in the image data, hence it can help to reduce the impact of variations in the lighting and color properties of the image.

% \begin{figure}[!h]
% 	 \centerline{\includegraphics[width=0.8\textwidth]{src/augmentations/Normalize.png}}
% 	 {\caption{Examples of Normalize Augmentation with the means respectively for three channels are 0.485, 0.456 and 0.406 and standard deviations are 0.229, 0.224 and 0.225.} \label{fig:Normalize}}
% \end{figure}

\subsubsection{ToGray} \label{ToGray}

One advantage of converting images into grayscale is that it can help to reduce the impact of color biases in the dataset. Color biases can occur when images of a particular class are predominantly of a particular color. By converting the images into grayscale, the model can be trained to recognize objects of a particular class based on their shapes and textures, rather than their color. However, the obvious disadvantage of converting images into grayscale is that it can result in the loss of important color information. Color Patterns may still be important for distinguishing different types of marine wildlife, and for identifying specific features of the environment that are important for object recognition. Figure \ref{fig:ToGray} illustrates our examples in grayscale.

\begin{figure}[!h]
	 \centerline{\includegraphics[width=0.8\textwidth]{src/augmentations/ToGray.png}}
	 {\caption{Examples of ToGray} \label{fig:ToGray}}
\end{figure}

\subsubsection{ToSepia} \label{ToSepia}

\begin{figure}
	 \centerline{\includegraphics[width=0.8\textwidth]{src/augmentations/ToSepia.png}}
	 {\caption{Examples of ToSepia} \label{fig:ToSepia}}
\end{figure}
Sepia is a type of image colorization technique that adds a brownish or yellowish tint to an image, giving it a vintage or aged look. It is often used in artistic or design applications to give images a nostalgic or classic feel.

While converting an image to grayscale  \ref{ToGray} involves discarding color information and representing the image in shades of gray, applying a sepia tint preserves the original color information and modifies it to create a vintage or aged effect. This makes the sepia technique a good option for applications such as edge detection, where the shape and structure of objects in the image are more important than their color, without the risk of significant loss of color information.

\subsubsection{Gaussian Noise} \label{GaussNoise}
\begin{figure}
	 \centerline{\includegraphics[width=0.8\textwidth]{src/augmentations/GaussNoise.png}}
	 {\caption{Examples of Gaussian Noise Injection with variance range for noise between 10 and 50} \label{fig:GaussNoise}}
\end{figure}

Gaussian noise injection is an image augmentation technique that adds random noise to the image by sampling from a normal distribution. In this technique, a random value is added to each pixel of the image, and the value of the added noise is determined by the standard deviation of the normal distribution, as illustrated in Figure \ref{fig:GaussNoise}.

The advantage of Gaussian noise injection is that it can create a diverse set of training images with varying levels of noise. This can help the machine learning model to better recognize objects in images with varying levels of noise, which can occur due to factors such as lighting, camera settings, or weather conditions. By training the model on a diverse set of noisy images, it can learn to better generalize to new, unseen data with varying levels of noise.

However, in the context of aerial images of marine wildlife, Gaussian noise injection may not be the most effective image augmentation technique. Aerial images of marine environments are often captured from a high altitude and are subject to atmospheric distortions, which can result in noise that is different from the type of noise added by Gaussian noise injection. Gaussian noise injection can negatively impact the visual quality of the image, which can make it more challenging to detect and recognize objects in the image.